\section{Current proof boundary after the component, period and local-functor joins}
\label{sec:dependency-audit-v57}

\paragraph{The purpose of this audit.}
We continue to accept the pinned formal proof as correct.  The question
here is which portions now have a readable mathematical proof, and which
original constructions still need one.  This edition adds @NEW@ claims
to the @OLD@ preserved claims, for @TOTAL@ in all.  Every new claim has
a separate review bound to its exact proof text and source statements.
The records distinguish a complete original implication from a construction
of all of that implication's inputs.

\paragraph{What has moved.}
The arithmetic chapter now reaches the full original rank-two component
theorem.  Its local chains, dimension inequalities, characteristic-prime
curves, homogeneous incidence and seed localization are proved on the
actual rings.  Its original finite-weight, automorphic propagation and
seed assumptions are retained.  The rank-three parity and exceptional
component arguments remain distinct.

The period chapter proves the original universal Gram expansion at its
precise analytic inputs.  It establishes the scalar absolute majorant
before changing sums and integrals, and retains the original rational
returns, measures, finite phase, column convention and normalizing scalar.
It does not infer a uniform spectral estimate from an operator identity.

The geometric chapters prove the full determinant-jet window and carry it
to actual integral row germs, with an explicit unit determinant factor.
The original Cartier normal form remains the input to that last
application.  They also remove the full original section-growth input:
one very ample power has exactly polynomial section dimension at every
positive multiple.  The classical theorems used for this replacement and
the dictionary identifying their objects are printed in the proof.

The local period-ring chapter constructs a genuinely fixed integral basis,
the nonlinear Frobenius conjugacy, and the original etale functor with
every lattice morphism.  The earlier linear solution-descent theorem
alone did not supply this conjugacy.  The new construction supplies the
actual etale input to the previously proved conditional Wach theorem.
Herr class maps, quantitative comparison and the common global family
remain to be reconstructed at their full scope.

\paragraph{Whole obligations and shared objects.}
The selected sufficient graph retains its 121 nodes and 161 edges.
There remain @OPEN@ whole selected obligations: @C@ class C and @B@
class B.  Separately, @SHARED@ are shared by RH and GRH and @EXTRA@
belong to the extra unrestricted Dirichlet-GRH route.  There is no new
whole selected closure.  The earlier full
\texttt{GLOBAL\_LOCAL\_GAIN} closure is preserved.

The unchanged whole count does not undo the completed subtheorems.
These large obligations combine several constructions and their
compatibility conditions.  A completed local ingredient becomes useful
only through its original maps on the common objects; its existence on
an independently chosen replacement object would not supply that join.
The register below records both the supplied ingredient and the exact
remaining construction, so subsequent work can begin at the first
unproved input rather than repeating earlier results.

\paragraph{Reading and review conventions.}
Read ACC for arithmetic components, PGR for periods, DJW for determinant
jets and integral row tests, SGV for the section-growth replacement,
and FET for faithful invariants and the local functor.  Published
foundations are explicit inputs, not unlabelled proof gaps.  The proof
and source hashes ensure that a review refers to the distributed text;
they are not a substitute for mathematical review or formal verification.
Draft TeX checks syntax and references.  No PDF rendering or new Lean
build is asserted.

\subsection{Current register of the complete remaining contracts}

